\section{Dictionaries (\texttt{dict}) --- Key-Value Associative Mappings}

A \texttt{dict} is Python's main structure for associating one object with another. Instead of asking for a component by numeric position, a dictionary asks for a value by key. The key acts as the lookup handle, and the value is the stored object attached to that key.

This section studies dictionaries as hash-based associative mappings. We will examine their key-value syntax, the requirement that keys be hashable and hash-stable, the freedom of dictionary values, construction patterns, insertion-order behavior, dynamic view objects, safe and unsafe mutation patterns, merging, lookup, and deletion.

The central idea is simple: a dictionary is not a list of pairs with nicer syntax. It is a runtime mapping engine where stable keys provide direct access to associated values.